% Section 3 -- Problem formulation (IEEEtran)
% Source of every claim: docs/registre_affirmations.md ([V] stated as fact, [H] stated as hypothesis, [Q] omitted).
% Each factual or numeric sentence is followed by a "% REG:" comment giving the registry row(s).
\section{Problem Formulation}
\label{sec:formulation}

This section states the decision problem that the learning experiments address.% STRUCT
The platform, the dataset and the canonical pipeline are described in Sec.~\ref{sec:data_pipeline}; the window construction and the cost definition are recalled here only as far as the formulation needs them.% REG: WIN-a, COST-a

\subsection{Scope: One-Step Configuration Selection}
\label{sec:scope}

The dataset contains no transition in which the configuration changes (Sec.~\ref{sec:config_constant}).% REG: D1-a, D1-c
It therefore cannot identify the effect of an action on the next state, and we do not model the state transition kernel.% REG: D1-c, NO-seq
We formulate the task as a \emph{constrained contextual bandit}: at each decision window the selector observes a context $\mathbf{x}$ and chooses one of eight configurations, and the discount factor of the problem is $\gamma=0$.% REG: D1-c, NO-seq
The choice $\gamma=0$ follows from D1 in Sec.~\ref{sec:config_constant}: with a configuration that is constant within every reservation, no logged data constrains a value that bootstraps over successive windows.% REG: D1-a, D1-c
It is a modelling choice and not a finding that the effect of an action is confined to the current window.% REG: D1-c, E2-c
The agents of the main comparison are trained with a discount of $\gamma=0.7$ on consecutive windows of an episode and are used as one-step selectors, each window receiving the greedy action of the learned network.% REG: E12P-b, E12P-d
The variant with $\gamma=0$, which matches the one-step objective, is an ablation, and its outcome depends on the set of scripts considered.% REG: E12-i, E12P-b

We do not claim sequential control of the slicing, nor that the Markov property holds for the state below.% REG: NO-seq, D1-c
Evaluation is at the reservation level, and sequential control is left as future work.% REG: NO-seq

We frame the selector as a component of a slower loop, in the style of an rApp hosted by the Non-RT RIC, rather than as an xApp of the Near-RT RIC~\cite{TODO-oran-timescales}.% REG: NO-nearrt
This framing is a choice that has to be justified, not a result: the decision window is 5\,s, the files do not allow the 10\,ms--1\,s range for either signal, and we run no window-sensitivity experiment (Sec.~\ref{sec:sampling}).% REG: WIN-a, E13b-d, E13b-e, NO-nearrt
We do not simulate a deployed control loop, since the dataset contains no action change to validate one against.% REG: D1-c, NO-nearrt

\subsection{Context}
\label{sec:state}

A decision instance is one non-overlapping window of 5\,s, with the window index defined in Sec.~\ref{sec:data_pipeline}.% REG: WIN-a
The context $\mathbf{x}_t\in\mathbb{R}^{10}$ has two groups of variables,
\begin{equation}
  \mathbf{x}_t=\bigl[\,n_t,\ \mathrm{rsrp}_t,\ \mathrm{snr}_t,\ \mathrm{mcs}_t,\ \mathrm{bler}_t,\ b_t,\ \rho_{t-1},\ \rho_{t-2},\ \ell_{t-1},\ \ell_{t-2}\,\bigr],
  \label{eq:state}
\end{equation}% REG: STATE-a
where $n_t$ is the number of active UEs; $\mathrm{rsrp}_t$, $\mathrm{snr}_t$, $\mathrm{mcs}_t$, $\mathrm{bler}_t$ and $b_t$ are the downlink radio indicators and the uplink buffer status (\texttt{ul\_buff}) of the current window, averaged over the samples of the attached eMBB UEs; and $\rho_{t-1},\rho_{t-2}$ and $\ell_{t-1},\ell_{t-2}$ are the one- and two-window lags of the eMBB reward and of the median URLLC latency.% REG: STATE-a, REW-a
The original notebook also placed the reward $\rho_t$ and the median latency $\ell_t$ of the current window in the state, which gave 12 dimensions; they are removed, and only their lags remain.% REG: AU-a
Lags are formed inside an episode, and a gap in the windows starts a new episode (Sec.~\ref{sec:data_pipeline}).% REG: D3-a

The configuration is not a variable of $\mathbf{x}_t$, but the context is not free of information about it.% REG: D4-a, NO-dyn
On the 12-dimensional state of the original notebook, the lagged variables alone predict the logged action with 30.5--32.8\% accuracy over the eight actions, against a chance level of 12.4\%.% REG: D4-a
These accuracies have not been re-measured on the 10-dimensional state above.% REG: D4-a, D4-d
Moreover, the lags are outcomes recorded under the logged configuration, which is constant within a reservation, so the context contains outcomes of the action that the selector is asked to choose.% REG: D1-a, D4-a
We therefore do not interpret the selector as having learned an effect of the configuration that is independent of the logged one, and a learned selector may partly act as a classifier of the logged configuration.% REG: NO-dyn, D4-a
We hypothesize that adding $a_{t-1}$ to the context would add little, because the action is constant within a reservation and is therefore already predictable; this is an argument and not an experimental result.% REG: D4-c

\subsection{Actions}
\label{sec:actions}

An action selects the PRB split between the two slices and the MAC scheduler,
\begin{equation}
  a=(p,m),\qquad p\in\{9,21,30,39\},\quad m\in\{\mathrm{RR},\mathrm{PF}\},
  \label{eq:action}
\end{equation}% REG: PLAT-a
where $p$ is the number of PRBs of the eMBB slice, which gives $|\mathcal{A}|=8$ actions (Table~\ref{tab:actions}).% REG: PLAT-a
The fifth slicing of the dataset, listed as 50/0 in the dataset README, is not an action; the README label and the field \texttt{slicing\_enabled=0} are not reconciled (Sec.~\ref{sec:data_pipeline}).% REG: D1-b, PLAT-a

\begin{table}[t]
\centering
\caption{Action set: four eMBB/URLLC PRB splits $\times$ two schedulers.}
\label{tab:actions}
\begin{tabular}{ccc}
\hline
eMBB/URLLC PRBs & Scheduler RR & Scheduler PF \\
\hline
9/41  & $(9,\mathrm{RR})$  & $(9,\mathrm{PF})$ \\
21/29 & $(21,\mathrm{RR})$ & $(21,\mathrm{PF})$ \\
30/20 & $(30,\mathrm{RR})$ & $(30,\mathrm{PF})$ \\
39/11 & $(39,\mathrm{RR})$ & $(39,\mathrm{PF})$ \\
\hline
\end{tabular}
\end{table}% REG: PLAT-a

All eight actions are logged, but their support is uneven.% REG: D6-c, D5-a
Only three of the 19 traffic scripts cover all eight actions with usable reservations, and the eight other scripts with usable reservations cover 1 to 5 of the eight configurations.% REG: D6-c
The cells of the cost table are well supplied with windows but poor in reservations.% REG: D5-a
The dataset records one configuration per reservation, so a decision per window is a modelling convention: each window is treated as an instance whose logged action is the configuration of its reservation.% REG: D1-a, D1-c

\subsection{Reward}
\label{sec:reward}

The reward $r_t$ is the throughput of the eMBB slice in Mbps over the window: for each attached eMBB UE we average the downlink bitrate over the samples of the window, and we sum these means over the UEs.% REG: REW-a, D2-b
The eMBB windows of the UEs share the common origin $t_0$ of Sec.~\ref{sec:clocks}; the original notebook summed windows that were shifted per UE.% REG: D2-b, WIN-a
The reward is read from the logs, with no performance model.% REG: AU-a
The estimated alignment of the UE clocks is an assumption that we have not validated independently (Sec.~\ref{sec:clocks}).% REG: D2-a, D2-b
With the slicing fixed, the difference in reward between PF and RR changes sign across the 12 (script, slicing) pairs and is not distinguishable from zero.% REG: D7-a

\subsection{Cost}
\label{sec:cost}

The cost $c_t\in[0,1]$ is the fraction of the URLLC packets of the window whose one-way latency exceeds $L_{\max}=20$\,ms,
\begin{equation}
  c_t=\frac{1}{N_t}\sum_{i=1}^{N_t}\mathbb{1}\bigl(\ell_i>L_{\max}\bigr),
  \label{eq:cost}
\end{equation}% REG: COST-a, E13d-a
where $\ell_i=(\text{Received Time}-\text{Sent Time})\times 10^{3}$\,ms, and the $N_t$ packets are those kept by the latency filter ($0\le\ell_i<20\,000$\,ms) and pooled over the readable URLLC UEs of the reservation.% REG: COST-a, E13d-a
A window without any kept packet has no cost row.% REG: COST-a
The filter is analysed in Sec.~\ref{sec:latency}; that analysis covers the training set only.% REG: E13d-h, NO-filter
The cost is a rate over a window and not a tail statistic; we report no analysis of the 99th-percentile latency.% REG: R1W2-no
The budget is $d=0.03$, and we do not document a source for $d$ or for $L_{\max}$~\cite{TODO-budget-latency-source}.% REG: COST-a, NO-source

The resolution of $c_t$ is limited.% REG: E13b-c
With 10 packets per second per UE and 2 to 3 active UEs, the criterion of at least 34 packets (one violation out of 34 is below $d$) in 99\% of the windows is met at 5\,s for only 7 of the 12 reservations examined.% REG: E13b-c

During training, the cost enters the constraint through a table $\hat c(a,\beta)$ of the mean window cost for each action $a$ and each load bin $\beta$.% REG: AU-b
The pipeline uses a single table with three bins of $n_t$, with edges $[0;3.5;5.5;8.1]$ and a fallback for cells with fewer than 20 windows, which was not triggered in the canonical pipeline.% REG: AU-b
We select no number of load levels as optimal (Sec.~\ref{sec:cells}).% REG: NO-bins, E13a2-c
The table is a coarse estimate: between 47\% and 63\% of its cells have a reservation-level confidence interval that contains $d$.% REG: D5-b
It also treats the load bin as if it were independent of the action, which the data do not show: the number of active UEs depends on the action across the eight actions.% REG: D7-c

\subsection{Constrained Objective}
\label{sec:objective}

Let $\pi(a\mid\mathbf{x})$ be a stochastic selector.% STRUCT
We seek
\begin{equation}
  \max_{\pi}\ J_R(\pi)=\mathbb{E}\bigl[r(\mathbf{x},a)\bigr]
  \quad\text{s.t.}\quad
  J_C(\pi)=\mathbb{E}\bigl[c(\mathbf{x},a)\bigr]\le d,\qquad d=0.03,
  \label{eq:objective}
\end{equation}% REG: COST-a, D1-c, NO-seq
where $\mathbf{x}$ is drawn from the distribution of contexts, $a\sim\pi(\cdot\mid\mathbf{x})$, and $r(\mathbf{x},a)$ and $c(\mathbf{x},a)$ denote the expected reward and the expected cost of action $a$ in context $\mathbf{x}$.% REG: D1-c
The constraint is a bound on the \emph{mean per-window cost}; it is not a bound on the cost of each reservation, and it carries no probabilistic guarantee.% REG: NO-safe, E13d-d
Under a relaxation with multiplier $\lambda\ge0$ the problem becomes the maximization of
\begin{equation}
  \mathcal{L}(\pi,\lambda)=J_R(\pi)-\lambda\bigl(J_C(\pi)-d\bigr),
  \label{eq:lagrangian}
\end{equation}% REG: NO-budget
a standard relaxation of constrained problems~\cite{TODO-lagrangian-relaxation}; the estimation of the cost and the update of $\lambda$ are described in Sec.~\ref{sec:method}.% REG: AU-b
The budget is feasible on the tables of the training scripts: the minimal attainable cost is 0.0211 for script a908e4cc, 0.0112 for script a53521b2 and 0.0182 for the whole training set, all read in tables.% REG: E1-a
Feasibility on tables does not carry over to training: no run of the diagnostic satisfies the frozen convergence criterion (0 of 60), and no value of $\alpha$ in the grid is feasible on validation.% REG: E2-a, E11-a
We therefore do not claim that the solution of Eq.~\eqref{eq:objective} is reached, or that the learned selector respects $d$.% REG: NO-budget, E12-e

\subsection{Assumptions and Non-Claims}
\label{sec:assumptions}

The formulation rests on the following assumptions; none is verified, and some are contradicted by the audit.% REG: D7-c, D7-b, D6-b
\begin{enumerate}
  \item[(A1)] \emph{Exogenous context.} The load and the radio context are not influenced by the action. The MGEN traffic is open loop (\texttt{UDP}, \texttt{PERIODIC}), but the emitted rate being identical across reservations is unverified, and the number of active UEs depends on the action across the eight actions (Sec.~\ref{sec:clocks}).% REG: D6-f, D7-c
  \item[(A2)] \emph{Configuration as good as assigned.} Given the context, the logged configuration is independent of unobserved factors. For the scheduler this is contradicted: PF and RR occupy disjoint blocks of reservations, 13\,h to about 8 days apart, so the scheduler is not separable from the execution time (Sec.~\ref{sec:confound}).% REG: D7-b, D8-b
  \item[(A3)] \emph{Independent unit.} The independent unit for generalization is the script; this hypothesis is strongly supported by the structure of the data but not tested (Sec.~\ref{sec:scripts}).% REG: D6-b
\end{enumerate}

With these assumptions, we make no claim of sequential control, of a Markov property, of a safety guarantee, of Near-RT validation, of generalization to other scripts, or of superiority of a learned selector over a table built on the same cells.% REG: NO-seq, NO-safe, NO-nearrt, NO-general, NO-rl
We describe the scheduler association as an association and not as a causal effect.% REG: NO-pfcause
